\documentclass[10pt,twocolumn,letterpaper]{article}

\usepackage[pagenumbers]{cvpr} % CVPR template formatting
\usepackage{graphicx}
\usepackage{amsmath}
\usepackage{amssymb}
\usepackage{booktabs}
\usepackage{multirow}
\usepackage{microtype}
\usepackage{xcolor}
\usepackage[pagebackref,breaklinks,colorlinks]{hyperref}

\def\paperID{P24}
\def\confName{CS-UCS532}
\def\confYear{2026}

\title{Scientific Disentanglement and Empirical Audit of Semi-Supervised Action Recognition: An Analysis of $F^2$-FixMatch on UCF-101}

\author{
Team P24\\
Department of Computer Science and Engineering\\
{\tt\small [GitHub: https://github.com/MankiratSingh-git/P24-CV]}
}

\begin{document}
\maketitle

\begin{abstract}
Semi-supervised action recognition seeks to leverage unannotated video to close the performance deficit toward full supervision. Recently, $F^2$-FixMatch reported a +21.41\% gain over standard single-stream FixMatch on the UCF-101 10\% labeled benchmark, attributing this improvement to a novel feature-level consistency module comprising Progressive Augmentations (P-Aug) and Multi-branch Pseudo-Labeling (MPL). In this work, we conduct a rigorous methodological audit and empirical disentanglement of this claim. By introducing the missing experimental control—a dual-stream RGB+Temporal Gradient (TG) FixMatch baseline without feature-level losses ($\beta=0, \gamma=0$)—we mathematically decompose the observed gain. Physical evaluation on an identical 30-view multi-crop protocol reveals that the bi-modal RGB+TG dual-stream configuration accounts for $\mathbf{+20.54}$ points (\textbf{95.9\%} of the total gain), whereas the feature-level module contributes an incremental $\mathbf{+0.87}$ points (\textbf{4.1\%}). Furthermore, we provide the first empirical ablation of the confidence threshold on video FixMatch ($\tau = 0.95$ vs. $\tau = 0.30$), proving that data starvation from strict thresholding degrades Top-1 accuracy by $-1.56\%$. Finally, we demonstrate that cross-modal BatchNorm contamination suppresses clean RGB inference, and recalibrating statistics via PreciseBN immediately recovers $+2.91\%$ without altering learned weights. All code, configs, and logs are publicly available.
\end{abstract}

\section{Introduction}
\label{sec:intro}
Deep video action recognition architectures typically require extensive supervised video annotations, which are prohibitively expensive to collect. Semi-supervised learning (SSL) provides an attractive paradigm: learning from a tiny fraction of annotated videos while exploiting large unlabeled video pools.

On the canonical UCF-101 benchmark with 10\% labeled training data (953 labeled and 8,584 unlabeled videos), standard single-stream FixMatch yields 38.49\% Top-1 accuracy. In contrast, the recently proposed $F^2$-FixMatch framework reaches 59.90\%, closing over 60\% of the deficit toward full supervision (72.91\%). This headline gain of $+21.41\%$ was attributed in the literature to feature-level consistency machinery (P-Aug and MPL). 

However, $F^2$-FixMatch differs from baseline FixMatch along multiple concurrent axes: input modalities (introducing Temporal Gradients), dual-head prediction architectures, pseudo-label ensembling, and multi-term objective functions. Prior to this study, no ablation isolated these intertwined factors.

In this work, we execute a rigorous, fail-stop scientific investigation to answer three fundamental questions:
\begin{enumerate}
    \item \textbf{Attribution}: What fraction of the $+21.41\%$ improvement stems from the bi-modal RGB+TG dual-stream pipeline versus the feature-level consistency module?
    \item \textbf{Threshold Dynamics}: Does the canonical 2D FixMatch threshold ($\tau = 0.95$) benefit video SSL, or does the lower threshold ($\tau = 0.30$) provide necessary coverage?
    \item \textbf{Normalization Drift}: Does sharing batch normalization parameters across RGB and Temporal Gradient streams corrupt evaluation statistics?
\end{enumerate}

\section{Related Work}
\label{sec:related}
\textbf{Consistency-Based Semi-Supervised Learning.} FixMatch~\cite{sohn2020fixmatch} established weak-to-strong consistency regularization with high-confidence pseudo-label thresholding ($\tau = 0.95$) as the dominant paradigm in 2D image recognition. Extensions such as FlexMatch~\cite{zhang2021flexmatch} and FreeMatch~\cite{wang2023freematch} dynamically adjust thresholds to account for class-dependent learning difficulties.

\textbf{Semi-Supervised Video Action Recognition.} Adapting SSL to 3D video backbones is challenging due to large memory footprints and temporal dynamics. The VideoSSL benchmark~\cite{xiao2021videossl} established standard UCF-101 splits and identified that 3D-ResNet-18 models require lower confidence thresholds ($\tau = 0.30$) due to higher entropy in spatio-temporal representations. Multi-view pseudo-labeling (MvPL)~\cite{kim2022mvpl} demonstrated that fusing complementary modalities, such as optical flow or temporal gradients, substantially strengthens pseudo-label quality.

\textbf{Feature Perturbations and Dual Streams.} $F^2$-FixMatch~\cite{tu2025f2fixmatch} combined frame-level FixMatch with feature-level progressive augmentations (P-Aug) and auxiliary Temporal Average Pooling (TAP) heads (MPL). However, as their official implementation code was unavailable, empirical verification of the relative importance of these components remained open.

\section{Methodology \& Implementation Audits}
\label{sec:method}
\subsection{System Architecture}
Our system processes 8-frame clips sampled with stride 8 ($224 \times 224$ pixels) through a 3D-ResNet-18 backbone (31.92M parameters). The framework ingests two streams:
\begin{enumerate}
    \item \textbf{RGB Stream ($x^V$)}: Raw video appearance.
    \item \textbf{Temporal Gradient ($x^G$)}: Consecutive frame differences, defined as $x_t^G = (x_t^V - x_{t+1}^V + 255) / 2$.
\end{enumerate}
A single shared backbone extracts representations for both streams. For unlabeled inputs, weak augmentations generate ensemble pseudo-labels:
\begin{equation}
    p_w = \frac{1}{2} \left( \text{Softmax}(C_{3D}(Z_w^{\text{RGB}})) + \text{Softmax}(C_{3D}(Z_w^{\text{TG}})) \right)
\end{equation}
The pseudo-label is retained if $\max(p_w) \ge \tau$. Strong augmented views are supervised via cross-entropy loss $\mathcal{L}_u^s$.

\subsection{Feature-Level Module Audit}
The feature-level consistency module introduces auxiliary losses $\beta \mathcal{L}_u^{w_p} + \gamma \mathcal{L}_u^{s_p}$ over TAP representations. Our tensor audit revealed two critical architectural properties:
\begin{enumerate}
    \item \textbf{Weak Feature Detach ($Z_w$)}: In the canonical code, weak features $Z_w$ are detached prior to TAP projection. Consequently, $\mathcal{L}_u^{w_p}$ does not backpropagate into the backbone; it updates only the auxiliary head $C_{\text{TAP}}$, which is discarded at test time.
    \item \textbf{Spatial Pooling Invariance}: In the original implementation, spatial flipping is immediately followed by spatial average pooling, rendering spatial flip perturbations mathematically invariant at the head input.
\end{enumerate}
Rather than silently modifying these design choices, we preserve them in our canonical runs and isolate their collective influence using our experimental control.

\begin{table*}[t]
\centering
\small
\caption{\textbf{The 5-Rung Experimental Ladder on UCF-101 Split 1.} All models evaluated under the matched 30-view multi-crop testing protocol (10 temporal clips $\times$ 3 spatial crops, $256 \times 256$, official test set $N = 3,783$ videos).}
\label{tab:main_ladder}
\begin{tabular}{@{}cclccccrr@{}}
\toprule
\textbf{Run} & \textbf{Modality} & \textbf{Loss Weights} & \textbf{Description} & \textbf{Best Top-1 (\%)} & \textbf{Final Top-1 (\%)} & \textbf{Top-5 (\%)} & \textbf{$\Delta$ vs. 10\% Sup} & \textbf{$\Delta$ vs. FixMatch} \\
\midrule
\textbf{A} & RGB & Supervised & 100\% Supervised Oracle & \textbf{72.91} & — & 90.06 & +41.59 & +34.42 \\
\textbf{B} & RGB & Supervised & 10\% Supervised Baseline & 31.32 & 30.80 & 54.11 & \textit{Reference} & -7.17 \\
\textbf{C} & RGB & $\lambda_u=2.0$ & FixMatch RGB ($\tau=0.30$) & 38.49 & 38.88 & 64.00 & +7.17 & \textit{Reference} \\
\textbf{D} & \textbf{RGB + TG} & $\beta=0, \gamma=0$ & \textbf{Dual-Stream Control} & \textbf{59.03} & \textbf{59.37} & \textbf{82.90} & \textbf{+27.71} & \textbf{+20.54} \\
\textbf{F} & \textbf{RGB + TG} & $\beta=2, \gamma=1$ & Full $F^2$-FixMatch & \textbf{59.90} & \textbf{60.35} & \textbf{83.32} & \textbf{+28.58} & \textbf{+21.41} \\
\bottomrule
\end{tabular}
\end{table*}

\section{Experiments and Disentanglement}
\label{sec:experiments}

\subsection{Experimental Setup}
All models share a 3D-ResNet-18 backbone and are trained from scratch on UCF-101 Split 1 (953 labeled, 8,584 unlabeled) using an NVIDIA RTX A5000 GPU (24GB VRAM). We use SGD with momentum 0.9, weight decay $10^{-4}$, initial learning rate 0.2, cosine scheduling with a 40-epoch warmup, labeled batch size 4, unlabeled batch size 16 ($\mu=4$), and 2 gradient accumulation steps. Evaluation uses the standard 30-view protocol (10 clips $\times$ 3 spatial crops at $256 \times 256$).

\subsection{Main Benchmark Ladder}
Table~\ref{tab:main_ladder} presents the complete 5-rung methodological ladder. Moving from 10\% supervision (31.32\%) to single-stream FixMatch (38.49\%) yields $+7.17\%$. Full $F^2$-FixMatch achieves 59.90\%, closing 68.7\% of the deficit to 100\% full supervision.

\subsection{Disentanglement Attribution: Run D Control}
To decouple the bi-modal input stream from the feature consistency module, we executed the critical missing control: \textbf{Run D} (RGB+TG FixMatch with $\beta=0, \gamma=0$). The total semi-supervised gain is decomposed as:
\begin{align}
    \Delta_{\text{total}} &= F - C = +21.41\text{ pts (100.0\%)} \\
    \Delta_{\text{dual}} &= D - C = 59.03 - 38.49 = \mathbf{+20.54\text{ pts (95.9\%)}} \\
    \Delta_{\text{feat}} &= F - D = 59.90 - 59.03 = \mathbf{+0.87\text{ pts (4.1\%)}}
\textbf{Outcome \& Causal Scope}: These physical measurements unambiguously validate \textbf{Case C}: transitioning to the integrated dual-stream framework as a package accounts for 95.9\% of the observed improvement over single-stream FixMatch ($+20.54$ pts), while the feature-level consistency module delivers an incremental $+0.87\%$. Consequently, attributing the $+21.41\%$ leap primarily to P-Aug and MPL was an empirical misattribution.

Importantly, we emphasize that Run~D introduces the dual-stream configuration as an integrated package (bi-modal inputs, shared backbone parameterization, dual-branch supervised losses, cross-modal ensemble pseudo-labeling, and joint optimization dynamics). Our experiment rigorously establishes that this integrated package—rather than the feature consistency module—drives the performance jump, while leaving the finer-grained internal decomposition among individual sub-components (such as TG inputs alone vs. pseudo-label ensembling) as an open question for future work.

\subsection{Confidence Threshold Ablation: $\tau = 0.95$ vs. $0.30$}
To resolve whether video SSL requires lower thresholds, we trained FixMatch RGB with the canonical image threshold $\tau = 0.95$:
\begin{itemize}
    \item FixMatch ($\tau = 0.95$): \textbf{36.93\%} (Final: 37.09\%, Top-5: 63.81\%).
    \item FixMatch ($\tau = 0.30$): \textbf{38.49\%} (Final: 38.88\%, Top-5: 64.00\%).
\end{itemize}
While $\tau = 0.95$ improves pseudo-label precision from 74.54\% to 94.28\%, it reduces unlabeled sample coverage from 94.98\% (8,153 videos) to 41.54\% (3,566 videos). This causes **data starvation**, which degrades performance by $-1.56\%$. This empirically validates why video SSL models require $\tau \approx 0.30$.

\subsection{Cross-Modal BatchNorm Contamination}
In dual-stream training, shared BatchNorm3d layers blend statistics across RGB and TG passes. Re-estimating population statistics on clean RGB clips via PreciseBN boosts $F^2$-FixMatch from \textbf{59.90\% to 62.81\% (+2.91 pts)} with zero parameter changes, confirming that cross-modal normalization drift was an active drag on clean RGB inference.

\section{Failure Analysis}
\label{sec:failure}
\textbf{Pseudo-Label Noise at $\tau = 0.30$.} Evaluating pseudo-labels on the 8,584 unlabeled samples reveals that at $\tau = 0.30$, 25.46\% (2,076 videos) of accepted pseudo-labels are incorrect. Despite this label noise, the network gains $+7.17\%$ over supervised training, confirming high tolerance to label noise under weak-strong consistency.

\textbf{Motion-Deprived Categories.} Analysis of per-class accuracy reveals that classes with minimal camera or body motion (e.g., \textit{Hammering}: 3.03\%, \textit{Shotput}: 2.17\%) suffer because Temporal Gradient values hover near zero, depriving the TG stream of discriminative signal.

\section{Changes from Proposal}
\label{sec:changes}
In accordance with course guidelines, we document our pre-declared gates and changes:
\begin{enumerate}
    \item \textbf{Decoupling the Dual-Stream Framework}: The original proposal attributed the $+21\%$ gain to P-Aug and MPL. Following critical review, we introduced Run D to isolate the dual-stream framework as an integrated package, revealing that 95.9\% of the observed gain over FixMatch RGB occurs upon transitioning to the dual-stream pipeline, while P-Aug and MPL contribute an incremental $+0.87\%$.
    \item \textbf{Pruning Speculative Per-Class Claims}: We eliminated unverified qualitative failure stories (e.g., ``confused with PoleVault'') and replaced them with rigorous pseudo-label precision curves and PreciseBN diagnostics.
    \item \textbf{Standardized 30-View Protocol}: We eliminated mixed-protocol comparisons by evaluating all 5 rungs under the identical 30-view multi-crop testing standard.
\end{enumerate}

\section{Limitations}
\label{sec:limitations}
First, while PreciseBN narrows the deficit to the published 71.1\% by $+2.91\%$, an $\sim 8$-point gap remains, likely reflecting unreleased hyperparameter details and distributed batch scaling. Second, due to GPU budget constraints ($\sim 100$ GPU hours total), runs were conducted under a single seed ($42$).

\section{Contribution Statement \& AI Declaration}
\textbf{Contributions}: The authors conceived the disentanglement strategy, executed all GPU training and evaluations on the RTX A5000, implemented the diagnostic test suites, and analyzed the empirical outcomes.
\textbf{AI Declaration}: Antigravity AI (DeepMind) was utilized as a coding assistant for writing diagnostic scripts, formatting reports, and conducting lint checks.

{\small
\bibliographystyle{ieee_fullname}
\begin{thebibliography}{10}\itemsep=-1pt

\bibitem{sohn2020fixmatch}
K.~Sohn, D.~Berthelot, N.~Carlini, Z.~Zhang, H.~Zhang, C.~A. Raffel, E.~D. Cubuk, A.~Kurakin, and C.-L. Li.
\newblock Fixmatch: Simplifying semi-supervised learning with consistency and confidence.
\newblock In {\em NeurIPS}, 2020.

\bibitem{xiao2021videossl}
F.~Xiao, Y.~J. Lee, K.~Grauman, J.~Malik, and C.~Feichtenhofer.
\newblock A video is worth more than 1,000 words: Semi-supervised video representation learning via visual and audio mappings.
\newblock In {\em WACV}, 2021.

\bibitem{tu2025f2fixmatch}
Z.~Tu, et~al.
\newblock Leveraging frame- and feature-level progressive augmentation for semi-supervised action recognition.
\newblock {\em ACM TOMM}, 2025.

\bibitem{kim2022mvpl}
J.~Kim, et~al.
\newblock Multi-view pseudo-labeling for semi-supervised action recognition.
\newblock In {\em CVPR}, 2022.

\bibitem{zhang2021flexmatch}
B.~Zhang, et~al.
\newblock Flexmatch: Boosting semi-supervised learning with curriculum pseudo-labeling.
\newblock In {\em NeurIPS}, 2021.

\bibitem{wang2023freematch}
Y.~Wang, et~al.
\newblock Freematch: Self-adaptive thresholding for semi-supervised learning.
\newblock In {\em ICLR}, 2023.

\bibitem{wu2018gn}
Y.~Wu and K.~He.
\newblock Group normalization.
\newblock In {\em ECCV}, 2018.

\bibitem{soomro2012ucf101}
K.~Soomro, A.~R. Zamir, and M.~Shah.
\newblock Ucf101: A dataset of 101 human actions classes from videos in the wild.
\newblock {\em CRCV-TR-12-01}, 2012.

\end{thebibliography}
}

\end{document}
